% Part: many-valued-logic
% Chapter: infinite-valued-logics
% Section: introduction

\documentclass[../../../include/open-logic-section]{subfiles}

\begin{document}

\olfileid{mvl}{inf}{int}

\olsection{پېژندنه}

د يوه ماتريس د رښتياوالي د قيمتونو شمېر لازماً متناهي نۀ وي۔
د نامتناهي ډېرو رښتياوالي قيمتونو د سټ دپاره يوه څرګنده ټاکنه
د~$0$ او~$1$ تر منځ د ناطقو عددونو سټ دے، يعنې~$V_\infty = [0,1] \cap \Rat$:
\begin{align*}
    V_\infty & = \Setabs{\frac{n}{m}}{n,m \in \Nat \text{ او } 0<m \text{ او } n\le m}.
\intertext{د رښتياوالي د دې نامتناهي سټ د کتنې په وخت کښې ډېر ځله
د لاندې فرعي سټونو کتل هم ګټور وي:}
V_m & = \Setabs{\frac{n}{m-1}}{n \in \Nat \text{ او } n<m}
\intertext{د بېلګې په توګه، $V_5$ د $5$ په برابر واټن پرتو رښتياوالي قيمتونو سټ دے:}
V_5 & = \{0, \frac{1}{4}, \frac{1}{2}, \frac{3}{4}, 1\}.
\end{align*}
د سرچينې يادونه: په لومړۍ ښودنه کښې اصلي متن صفر مخرج نۀ منع کوي؛
دلته د مخرج مثبت شرط ورزيات شو۔ په دويمه ښودنه کښې د سرچينې حد
يو اضافي، له يو څخه لوے قيمت جوړوي؛ حد دلته سم شو۔ دا دويمه بڼه
يوازې د دوو يا تر دوو زياتو قيمتونو دپاره کارېږي۔
د دغو رښتياوالي قيمتونو پر سټونو ولاړ منطقونو کښې عموماً يوازې~$1$
ټاکل شوے وي؛ يعنې~$V^+ = \{1\}$۔ په بله وينا،~$1$ د (مطلق)
رښتيا او~$0$ د مطلق دروغ نقش لوبوي، خو !!{formula}s په~$V$
کښې هر منځنے قيمت هم اخيستلے شي۔

سره له دې، د~$0$ او~$1$ تر منځ د ټولو \emph{حقيقي} عددونو
سټ~$V_{[0,1]} = [0,1]$، يا د~$[0,1]$ نور نامتناهي فرعي سټونه
هم کارولے شو۔ په دې رښتياوالي سټ ولاړ منطقونه ډېر ځله \emph{فازي}
بلل کېږي۔


\end{document}
